\documentclass[11pt,a4paper]{article}
\usepackage[utf8]{inputenc}
\usepackage[margin=1in]{geometry}
\usepackage{amsmath,amssymb,amsthm}
\usepackage{listings}
\usepackage{xcolor}
\usepackage{hyperref}

\newtheorem{theorem}{Theorem}
\newtheorem{lemma}[theorem]{Lemma}

\lstset{
  language=C++,
  basicstyle=\ttfamily\small,
  keywordstyle=\color{blue}\bfseries,
  commentstyle=\color{green!60!black},
  stringstyle=\color{red!70!black},
  numbers=left,
  numberstyle=\tiny\color{gray},
  breaklines=true,
  frame=single,
  tabsize=4,
  showstringspaces=false
}

\title{IOI 2009 -- Hiring}
\author{Solution}
\date{}

\begin{document}
\maketitle

\section{Problem Statement}
There are $N$ candidates, each with minimum salary $S_i$ and qualification $Q_i$. The \emph{fairness} constraint requires a constant $c$ such that every hired worker~$i$ receives salary $c \cdot Q_i \ge S_i$. Given budget $W$, maximize the number of hired workers. Among solutions with the same count, minimize total cost.

\section{Solution}

\subsection{Key Observation}
Define the \emph{ratio} $r_i = S_i / Q_i$. Worker~$i$ is eligible when $c \ge r_i$. If we fix $c = r_k$ for some worker~$k$, all workers with $r_i \le r_k$ are eligible, and the total cost is $c \cdot \sum Q_i = r_k \cdot \sum Q_i$. To minimize cost (and thus hire more workers), we should select the eligible workers with the smallest $Q_i$ values.

\subsection{Algorithm}
\begin{enumerate}
    \item Sort workers by ratio $r_i$ in increasing order.
    \item Process workers in sorted order. For worker~$k$, set $c = r_k$.
    \item Maintain a max-heap of $Q$ values and a running sum of selected $Q$'s.
    \item Insert $Q_k$. While $c \cdot \text{sumQ} > W$ (equivalently $S_k \cdot \text{sumQ} > W \cdot Q_k$), remove the largest $Q$ from the heap.
    \item Update the best answer.
\end{enumerate}

\subsection{Complexity}
\begin{itemize}
    \item \textbf{Time:} $O(N \log N)$.
    \item \textbf{Space:} $O(N)$.
\end{itemize}

\section{C++ Solution}

\begin{lstlisting}
#include <bits/stdc++.h>
using namespace std;

int main(){
    ios::sync_with_stdio(false);
    cin.tie(nullptr);

    int N;
    long long W;
    cin >> N >> W;

    vector<long long> S(N), Q(N);
    vector<int> idx(N);
    for(int i = 0; i < N; i++){
        cin >> S[i] >> Q[i];
        idx[i] = i;
    }

    // Sort by ratio S[i]/Q[i] increasing (cross-multiply to avoid floats)
    sort(idx.begin(), idx.end(), [&](int a, int b){
        return S[a] * Q[b] < S[b] * Q[a];
    });

    priority_queue<long long> pq;
    long long sumQ = 0;
    int bestCount = 0;
    double bestCost = 1e18;
    int bestIdx = -1;

    for(int k = 0; k < N; k++){
        int i = idx[k];
        pq.push(Q[i]);
        sumQ += Q[i];

        // Budget check: S[i] * sumQ <= W * Q[i]
        while(!pq.empty() && (__int128)S[i] * sumQ > (__int128)W * Q[i]){
            sumQ -= pq.top();
            pq.pop();
        }

        int cnt = (int)pq.size();
        double cost = (double)S[i] / Q[i] * sumQ;
        if(cnt > bestCount || (cnt == bestCount && cost < bestCost)){
            bestCount = cnt;
            bestCost = cost;
            bestIdx = k;
        }
    }

    cout << bestCount << "\n";

    // Reconstruct the selected workers
    priority_queue<pair<long long,int>> pq2;
    sumQ = 0;
    for(int k = 0; k <= bestIdx; k++){
        int i = idx[k];
        pq2.push({Q[i], i});
        sumQ += Q[i];
        while(!pq2.empty() && (__int128)S[idx[bestIdx]] * sumQ > (__int128)W * Q[idx[bestIdx]]){
            sumQ -= pq2.top().first;
            pq2.pop();
        }
    }

    vector<int> selected;
    while(!pq2.empty()){
        selected.push_back(pq2.top().second + 1);
        pq2.pop();
    }
    sort(selected.begin(), selected.end());
    for(int x : selected) cout << x << "\n";

    return 0;
}
\end{lstlisting}

\section{Notes}
The greedy approach works because sorting by ratio $r_i$ ensures that as $c$ increases, new workers become eligible but the per-unit cost also increases. The max-heap efficiently evicts the most expensive workers (largest $Q$) to stay within budget. The use of \texttt{\_\_int128} avoids overflow when comparing $S_i \cdot \text{sumQ}$ against $W \cdot Q_i$.

\end{document}
